\documentclass[11pt]{amsart}
\usepackage{amsmath,amssymb,amsthm,booktabs}
\usepackage[margin=1.15in]{geometry}
\usepackage[hidelinks]{hyperref}
\newtheorem{theorem}{Theorem}[section]
\newtheorem{corollary}[theorem]{Corollary}
\newtheorem{proposition}[theorem]{Proposition}
\newtheorem{lemma}[theorem]{Lemma}
\theoremstyle{remark}
\newtheorem{remark}[theorem]{Remark}
\newcommand{\E}{\mathbb{E}}
\newcommand{\Prob}{\mathbb{P}}
\newcommand{\F}{\mathcal{F}}
\newcommand{\doi}[1]{\href{https://doi.org/#1}{doi:\nolinkurl{#1}}}
\newcommand{\arxiv}[1]{\href{https://arxiv.org/abs/#1}{arXiv:\nolinkurl{#1}}}
\setlength{\emergencystretch}{2.5em}

\begin{document}

\title[Forward-loop-erased ants at distance two]{Convergence of the forward-loop-erased ants process at distance two}
\author{Yitzchak Shmalo}
\address{Einstein Institute of Mathematics, The Hebrew University of Jerusalem, Jerusalem, Israel}
\email{yspennstate@gmail.com}
\subjclass[2020]{60K35, 60J10, 05C81}
\keywords{reinforced random walk, ants process, loop-erased random walk, P\'olya urn, shortest paths}

\begin{abstract}
In the forward-loop-erased ants process, the $n$th ant walks from a nest $N$ to a food source $F$ with transition
probabilities proportional to the current edge weights, and every edge of the chronological loop erasure of its path
gains weight one. Armfield and Mailler showed that when $N$ and $F$ are at distance two the ants find the shortest
paths, and left open whether $W(n)/n$ converges. We show that it converges almost surely and that every path of length
two keeps a positive share in the limit. The key fact is that for a common neighbour $i$ of $N$ and $F$ the relative
drift of $W_{\{N,i\}}(n)W_{\{i,F\}}(n)$ is the inverse effective conductance between $N$ and $F$, the same for all $i$,
up to a summable error; hence ratios of these products converge on every finite graph.
\end{abstract}

\maketitle

\section{Introduction}\label{sec:intro}

Kious, Mailler and Schapira~\cite{KMS22a,KMS22b} introduced a family of reinforced random walk models for the foraging
of ants. The edges of a finite graph $G=(V,E)$ carry weights $W_e(n)$, all equal to $1$ at time $0$. At step $n$ an ant
starts at the nest $N$ and performs a random walk on $G$ with conductances $W(n)$, stopped when it reaches the food
source $F$; a set of edges extracted from its trajectory is then reinforced, every edge of the set gaining weight one.
Different extraction rules give different models: the loop-erased and geodesic rules select a simple path from $N$ to
$F$, while the trace rule selects every edge the ant traversed. Kious, Mailler and Schapira conjectured that in the
backward-loop-erased version (the loop erasure of the time-reversed trajectory) and the uniform-geodesic version the
ants find the shortest paths and the process converges. They proved this for series-parallel graphs with source $N$ and
sink $F$ and for the lozenge graph respectively~\cite{KMS22a}. In the trace-reinforced version the ants do not always
find the shortest paths~\cite{KMS22b}. Mailler and Varin~\cite{MV26} proved convergence
for a two-nest process with the backward erasure on triangle-series-parallel graphs.

Armfield and Mailler~\cite{AM} studied the forward-loop-erased version, in which the reinforced path $\gamma_n$ is the
chronological loop erasure of the $n$th trajectory; the model and their distance-two theorem also appear in Armfield's
thesis~\cite[Chapter~2]{Arm26}. Their Theorem~1.1 states that if $\{N,F\}\notin E$ and $N$ and $F$
have $k\ge1$ common neighbours $1,\dots,k$, so that $S=\{\{N,i\},\{i,F\}:i\in[k]\}$ is the set of edges on shortest
paths, then almost surely
\begin{equation}\label{eq:AM}
\text{(a)}\ \ \frac{W_{\{N,i\}}(n)-W_{\{i,F\}}(n)}{n}\to0\quad(i\in[k]),\qquad
\text{(b)}\ \ \frac{W_e(n)}{n}\to0\quad(e\notin S).
\end{equation}
Other edges between the vertices $i\in[k]$ and the rest of the graph are allowed. They remark that convergence of
$W(n)/n$ is open; on the graph reduced to $S$ the path counts form a $k$-colour P\'olya urn, and they suggest that a
proof would need the weights off $S$ to vanish fast enough for that comparison to hold, as in the method
of~\cite{KMS22a}.

The main result answers this question without any rate for the weights off $S$.

\begin{theorem}\label{thm:main}
Under the assumptions of~\cite[Theorem~1.1]{AM} there is a random vector $(\lambda_1,\dots,\lambda_k)$ with
$\lambda_i>0$ for every $i$ and $\sum_i\lambda_i=1$ such that, almost surely,
\[
\frac{W_{\{N,i\}}(n)}{n}\to\lambda_i,\qquad \frac{W_{\{i,F\}}(n)}{n}\to\lambda_i\quad(i\in[k]),\qquad
\frac{W_e(n)}{n}\to0\quad(e\notin S).
\]
In particular the forward-loop-erased ants process converges, and every shortest path keeps a positive share of the
reinforcement in the limit.
\end{theorem}

The input from~\cite{AM} is their Theorem~1.1 and two identities for the probability that an edge at $N$ or at $F$ is
reinforced (\cite[Lemmas~2.12 and~2.13]{AM}), which we reprove in Section~\ref{sec:edges} for completeness. The new
ingredient is Proposition~\ref{prop:ratio}: for vertices $i,j$ adjacent to both $N$ and $F$, on any finite graph, the
ratio of the products $Q_i(n)=W_{\{N,i\}}(n)W_{\{i,F\}}(n)$ and $Q_j(n)$ converges almost surely to a finite positive
limit. The point is that in the conditional drift of $\log Q_i$ the probability $h_i$ that a walk from $i$ reaches $F$
before $N$ cancels: the edge at $N$ is reinforced with probability proportional to $W_{\{N,i\}}h_i$ and the edge at $F$
with probability proportional to $W_{\{i,F\}}(1-h_i)$, so the relative drift of the product is $1/C(n)$ up to a
summable error, the same for all $i$. Armfield and Mailler used a suitably normalized sum of these products as a
stochastic approximation in their proof of~\eqref{eq:AM}(a); here each $Q_i$ is followed separately. The square root
that turns Proposition~\ref{prop:ratio} into Theorem~\ref{thm:main} comes from~\eqref{eq:AM}(a), which gives
$Q_i(n)/n^2-(W_{\{N,i\}}(n)/n)^2\to0$.

We do not determine the law of $\lambda$. In simulations (Section~\ref{sec:numerics}) the log ratios of different
seeds are separated at $n=10^5$ and appear to stabilize, consistent with the P\'olya urn heuristic of~\cite{AM}: on the graph reduced to $S$, started
from unit weights, the
limit would have the uniform law on the simplex.

\section{Edge probabilities at the nest and at the food}\label{sec:edges}

Fix conductances $w=(w_e)_{e\in E}$ with $w_e>0$, write $\pi_v=\sum_{e\ni v}w_e$, let $X$ be the random walk with these
conductances started at $N$ and stopped at $\tau_F$, the hitting time of $F$, and let $\gamma=\mathrm{LE}(X)$ be its
chronological loop erasure. For $v\in V$ let $h(v)=\Prob_v(\tau_F<\tau_N)$ and let
\[
C=\pi_N\Prob_N(\tau_F<\tau_N^+)=\sum_{u\sim N}w_{\{N,u\}}h(u)
\]
be the effective conductance between $N$ and $F$.

\begin{lemma}[{\cite[Lemmas~2.12 and~2.13]{AM}}]\label{lem:edges}
For $u\sim N$, $\Prob(\{N,u\}\in\gamma)=w_{\{N,u\}}h(u)/C$. For $u\sim F$ with $u\ne N$,
$\Prob(\{u,F\}\in\gamma)=w_{\{u,F\}}(1-h(u))/C$.
\end{lemma}

\begin{proof}
The path $\gamma$ is simple, starts at $N$ and ends at $F$, so an edge at $N$ lies on $\gamma$ only as its first edge
and an edge at $F$ only as its last edge. During chronological loop erasure the current erased path always ends at the
current position of the walk, and a vertex is removed only when the walk returns to a vertex that precedes it on the
current path.

Let $\sigma$ be the time of the last visit of $X$ to $N$ before $\tau_F$. After $\sigma$ the walk does not return to $N$,
so $X_{\sigma+1}$ is never removed and the first edge of $\gamma$ is $\{N,X_{\sigma+1}\}$. The expected number of visits
to $N$ before $\tau_F$ is $1/\Prob_N(\tau_F<\tau_N^+)=\pi_N/C$, and each visit is the last one followed by a step to $u$
with probability $(w_{\{N,u\}}/\pi_N)\,h(u)$. Hence $\Prob(X_{\sigma+1}=u)=w_{\{N,u\}}h(u)/C$.

The vertex $F$ is first visited at time $\tau_F$, so the last edge of $\gamma$ is $\{X_{\tau_F-1},F\}$. Let $g(x,y)$ be
the expected number of visits to $y$ before $\tau_F$ by the walk started at $x$. Then
$\Prob(X_{\tau_F-1}=u)=g(N,u)\,w_{\{u,F\}}/\pi_u$. By reversibility $g(x,y)/\pi_y$ is symmetric in $x$ and
$y$, also for the walk killed at $F$~\cite[Exercise~2.1(e) and the proof of Proposition~2.1]{LP17}, and by the strong Markov property $g(u,N)=\Prob_u(\tau_N<\tau_F)\,g(N,N)=(1-h(u))\pi_N/C$.
Therefore $g(N,u)/\pi_u=g(u,N)/\pi_N=(1-h(u))/C$.
\end{proof}

\section{Ratios of products}\label{sec:ratio}

Let $(W(n))_{n\ge0}$ be the forward-loop-erased ants process on a finite connected graph $G$ with distinct vertices
$N,F$, let $\F_n$ be the $\sigma$-field generated by $W(0),\dots,W(n)$, and write $C(n)$ and $h_v(n)$ for the
quantities of Section~\ref{sec:edges} computed with $w=W(n)$. For a vertex $i\notin\{N,F\}$ adjacent to both $N$ and $F$
put
\[
A_i(n)=W_{\{N,i\}}(n),\qquad B_i(n)=W_{\{i,F\}}(n),\qquad Q_i(n)=A_i(n)B_i(n).
\]

\begin{proposition}\label{prop:ratio}
Let $i,j\notin\{N,F\}$ be vertices adjacent to both $N$ and $F$. Then $\log\bigl(Q_i(n)/Q_j(n)\bigr)$ converges almost
surely to a finite random variable.
\end{proposition}

\begin{proof}
Write $a_n=A_i(n+1)-A_i(n)$ and $b_n=B_i(n+1)-B_i(n)$, both in $\{0,1\}$, and
\[
\xi_n=\frac{Q_i(n+1)}{Q_i(n)}-1=\frac{a_n}{A_i(n)}+\frac{b_n}{B_i(n)}+\frac{a_nb_n}{Q_i(n)}.
\]
By Lemma~\ref{lem:edges}, $\E[a_n\mid\F_n]=A_i(n)h_i(n)/C(n)$ and $\E[b_n\mid\F_n]=B_i(n)(1-h_i(n))/C(n)$, so
\begin{equation}\label{eq:drift}
\E[\xi_n\mid\F_n]=\frac{1}{C(n)}+\frac{g_n}{Q_i(n)},\qquad g_n=\E[a_nb_n\mid\F_n]=\Prob\bigl(\gamma_{n+1}=(N,i,F)\mid\F_n\bigr).
\end{equation}
Since $\log(1+x)=x-r(x)$ with $0\le r(x)\le x^2/2$ for $x\ge0$, and $Q_i(0)=1$,
\begin{equation}\label{eq:decomp}
\log Q_i(n)=\sum_{m<n}\frac{1}{C(m)}+D_i(n)+M_i(n)-R_i(n),
\end{equation}
where $D_i(n)=\sum_{m<n}g_m/Q_i(m)$, $M_i(n)=\sum_{m<n}(\xi_m-\E[\xi_m\mid\F_m])$ and $R_i(n)=\sum_{m<n}r(\xi_m)$.

Each of the last three terms converges almost surely, by the following deterministic bound. If $a_m=1$ for the $\ell$th
time, then $A_i(m)=\ell$; hence $\sum_m a_m/A_i(m)^2\le\sum_{\ell\ge1}\ell^{-2}=\pi^2/6$, and the same holds for $b$ and
$B_i$. As $(a_mb_m/Q_i(m))^2\le a_m/A_i(m)^2$, we get $\xi_m^2\le3\bigl(2a_m/A_i(m)^2+b_m/B_i(m)^2\bigr)$ and
\[
\sum_{m\ge0}\xi_m^2\le\frac{3\pi^2}{2}.
\]
Thus $R_i(n)$ is nondecreasing and bounded by $3\pi^2/4$, and $M_i$ is a martingale with
$\E[M_i(n)^2]\le\E\sum_m\xi_m^2\le3\pi^2/2$, so it converges almost surely. Finally $D_i(n)$ is nondecreasing, and if
$a_mb_m=1$ for the $\ell$th time then $A_i(m)\ge\ell$ and $B_i(m)\ge\ell$, so
\[
\E\Bigl[\sum_{m\ge0}\frac{g_m}{Q_i(m)}\Bigr]=\E\Bigl[\sum_{m\ge0}\frac{a_mb_m}{Q_i(m)}\Bigr]\le\sum_{\ell\ge1}\frac{1}{\ell^2}<\infty,
\]
and $D_i(n)$ converges almost surely. The first sum in~\eqref{eq:decomp} does not depend on $i$, so subtracting the
decomposition for $j$ gives
$\log(Q_i(n)/Q_j(n))=(D_i-D_j)(n)+(M_i-M_j)(n)-(R_i-R_j)(n)$, which converges almost surely.
\end{proof}

\begin{remark}
The proposition uses neither the distance between $N$ and $F$ nor any lower bound on the hitting probabilities $h_i$.
On its own it does not say that $Q_i(n)$ grows like $n^2$; this is where~\eqref{eq:AM} enters.
\end{remark}

\section{Proof of Theorem~\ref{thm:main}}\label{sec:proof}

Every reinforced path is simple and starts at $N$, so exactly one edge at $N$ gains weight at each step and
$\sum_{e\ni N}W_e(n)=n+\deg N$. By~\eqref{eq:AM}(b) the edges at $N$ outside $S$ have $W_e(n)/n\to0$, hence with
$s_i(n)=A_i(n)/n$ and $t_i(n)=B_i(n)/n$,
\[
\sum_{i\in[k]}s_i(n)\to1,\qquad s_i(n)-t_i(n)\to0\quad(i\in[k])
\]
almost surely, the second by~\eqref{eq:AM}(a). The normalization $n+n_e$ of~\cite{AM} does not affect these limits,
since $|W_e(n)/n-W_e(n)/(n+n_e)|\le n_e(n+1)/(n(n+n_e))\to0$.
Since $s_i(n)\le(n+\deg N)/n\le1+\deg N$, it follows that $Q_i(n)/n^2-s_i(n)^2=s_i(n)(t_i(n)-s_i(n))\to0$.

By Proposition~\ref{prop:ratio}, $Q_i(n)/Q_1(n)\to\rho_i\in(0,\infty)$ almost surely, with $\rho_1=1$. As $Q_1(n)/n^2$ is
bounded,
\[
\frac{Q_i(n)}{n^2}-\rho_i\frac{Q_1(n)}{n^2}=\Bigl(\frac{Q_i(n)}{Q_1(n)}-\rho_i\Bigr)\frac{Q_1(n)}{n^2}\to0,
\]
and therefore $s_i(n)^2-\rho_is_1(n)^2\to0$. Using $|\sqrt x-\sqrt y|\le\sqrt{|x-y|}$ we get
$s_i(n)-\sqrt{\rho_i}\,s_1(n)\to0$ for every $i$. Summing over $i$ gives $s_1(n)\sum_l\sqrt{\rho_l}\to1$, so
\[
s_i(n)\to\lambda_i:=\frac{\sqrt{\rho_i}}{\sum_{l\in[k]}\sqrt{\rho_l}}\in(0,1],\qquad t_i(n)\to\lambda_i,
\]
and $\sum_i\lambda_i=1$, with $\lambda_i<1$ when $k\ge2$. Together with~\eqref{eq:AM}(b) this is
Theorem~\ref{thm:main}. Finally, a reinforced path other than $(N,i,F)$ that uses the edge $\{N,i\}$ also uses an edge
outside $S$, so the number of such paths up to time $n$ is $o(n)$ by~\eqref{eq:AM}(b), and the frequency of the path
$(N,i,F)$ among the first $n$ reinforced paths also tends to $\lambda_i$. \qed

\section{Simulations}\label{sec:numerics}

The code in the repository~\cite{Repo} checks Lemma~\ref{lem:edges} against $2\cdot10^5$ simulated loop-erased paths
for random conductances on two graphs (all deviations are within $1.4$ standard errors), and runs the process from unit
weights to $n=10^5$ on both graphs with three seeds each. Graph~A has common neighbours $2,3$ of $N=0$ and $F=1$, extra
vertices $4,5$ with edges $\{0,4\},\{4,2\},\{3,5\},\{5,1\}$ and chords $\{2,3\},\{4,5\}$; graph~B has common neighbours
$2,3,4$, extra vertices $5,6$ with edges $\{0,5\},\{5,2\},\{5,6\},\{6,4\},\{6,1\}$ and chords $\{2,3\},\{3,4\}$. The
table lists $\log(Q_3(n)/Q_2(n))$.

\begin{center}
\begin{tabular}{lrrrrr}
\toprule
 & $n=2\cdot10^3$ & $10^4$ & $3\cdot10^4$ & $6\cdot10^4$ & $10^5$\\\midrule
A, seed 1 & $-3.915$ & $-3.955$ & $-3.905$ & $-3.881$ & $-3.865$\\
A, seed 2 & $-3.530$ & $-3.544$ & $-3.478$ & $-3.495$ & $-3.490$\\
A, seed 3 & $-2.390$ & $-2.302$ & $-2.265$ & $-2.280$ & $-2.293$\\
B, seed 1 & $-0.643$ & $-0.737$ & $-0.745$ & $-0.752$ & $-0.731$\\
B, seed 2 & $-0.338$ & $-0.415$ & $-0.364$ & $-0.334$ & $-0.333$\\
B, seed 3 & $0.533$ & $0.523$ & $0.515$ & $0.550$ & $0.540$\\
\bottomrule
\end{tabular}
\end{center}

The log ratios settle early while the largest normalized weight off $S$ is still between $0.03$ and $0.07$ at $n=10^5$,
in line with the fact that Proposition~\ref{prop:ratio} needs no rate for the weights off $S$. The values for different
seeds are separated at $n=10^5$ and appear to stabilize; this is a finite-time observation and says nothing definite
about the law of the limit.

\section*{Acknowledgements and funding}
This work was done by multiple AI agents, Claude (Anthropic) and Codex (OpenAI), connected to
Math Brain~\cite{MathBrainRepo} and working together inside the Brain framework~\cite{BrainFramework}, a shared,
persistent memory in which the agents record claims,
proofs, failed attempts and methods, retrieve one another's results, coordinate who works on what, and audit each
other's work. For more on the Brain and Math Brain, see the Math Brain paper~\cite{MathBrainPaper}, the Math Brain
repository~\cite{MathBrainRepo} and the Brain framework repository~\cite{BrainFramework}.

The research presented in this paper was supported by the European Research Council (ERC) under the European Union's
Horizon Europe research and innovation programme (grant agreement No.~101041711), by the Simons Foundation as part of
the Collaboration on the Mathematical and Scientific Foundations of Deep Learning, by Heights Labs, by Convex Nexus
Capital, by the Israel Science Foundation (grant number 2258/19), and by the Israel Science Foundation (ISF Grant
4101/25). Computational resources were provided in part by Google Cloud through its education credits programme
(credit 530550103).

\section*{Declaration of competing interest}
The author holds equity in Heights Labs and in Convex Nexus Capital and has a management role at Convex Nexus Capital.
These entities are included in the funding declaration above.

\section*{Data availability}
The code that checks Lemma~\ref{lem:edges} and produces the table of Section~\ref{sec:numerics} is in the
repository~\cite{Repo}.

\section*{Declaration of generative AI and AI-assisted technologies}
This paper is part of a broader research project on AI proving problems. Claude agents (Anthropic) proved
Proposition~\ref{prop:ratio} and Theorem~\ref{thm:main} and prepared the manuscript; Codex agents (OpenAI) refereed the
proofs. The model is that of Kious, Mailler and Schapira~\cite{KMS22a} with the forward loop erasure of Armfield and
Mailler~\cite{AM}, whose Theorem~1.1 and edge identities are used here; the question of convergence is theirs.
Proposition~\ref{prop:ratio} and Theorem~\ref{thm:main} are the new results. The author directed the work and is
responsible for the paper.

\begin{thebibliography}{99}
\bibitem{AM} W.~Armfield and C.~Mailler, \emph{The forward-loop-erased ants walk}, preprint, 2026,
\arxiv{2610.18453}.
\bibitem{Arm26} W.~Armfield, \emph{An analysis of learning processes through reinforced random walks}, PhD thesis,
University of Bath, 2026.
\bibitem{KMS22a} D.~Kious, C.~Mailler and B.~Schapira, \emph{Finding geodesics on graphs using reinforcement learning},
Ann. Appl. Probab. \textbf{32} (2022), 3889--3929, \arxiv{2010.04820}.
\bibitem{KMS22b} D.~Kious, C.~Mailler and B.~Schapira, \emph{The trace-reinforced ants process does not find shortest
paths}, J. \'Ec. polytech. Math. \textbf{9} (2022), 505--536, \arxiv{2106.10559}.
\bibitem{LP17} R.~Lyons and Y.~Peres, \emph{Probability on Trees and Networks}, Cambridge Series in Statistical and
Probabilistic Mathematics~42, Cambridge University Press, 2016.
\bibitem{MV26} C.~Mailler and P.~Varin, \emph{The two-nest ants process on triangle-series-parallel graphs}, preprint,
2026, \arxiv{2601.22855}.
\bibitem{Repo} Y.~Shmalo, \emph{Forward-loop-erased ants at distance two: paper and simulation code}, 2026,
\url{https://github.com/yspennstate/forward-loop-erased-ants-convergence}.
\bibitem{MathBrainRepo} Y.~Shmalo, \emph{Math Brain: public snapshot, code and tests}, 2026,
\url{https://github.com/yspennstate/math-brain-public/tree/8c9a6bd478c4ec4082347151b7ae16db50c917dd}.
\bibitem{MathBrainPaper} Y.~Shmalo, \emph{Math Brain: shared, auditable memory for multi-agent mathematical research},
2026, \doi{10.5281/zenodo.24150837},
\url{https://github.com/yspennstate/math-brain-public/tree/main/papers/math-brain-2026-10-24}.
\bibitem{BrainFramework} Y.~Shmalo, \emph{Brain Framework}, public software repository, 2026,
\url{https://github.com/yspennstate/brain-framework-public/tree/5629e644cdd510173ae97e9aa92642dabd22dc57}.
\end{thebibliography}

\end{document}
